%%%%%%%%%%%%%%%%%%%%%%%%%%%%%%%%%%%%%%%%%%%%%%%%%%%%%%%%%%%%%%%%%%%%%%
% Lớp 10 | Thống kê | Số gần đúng. Sai số
% Dạng 0D6?1-1: Tính và ước lượng sai số tuyệt đối, tương đối
% Số câu: 14  (sinh tự động bởi tach_cau.py; từ nay thầy sửa trực tiếp file này)
%%%%%%%%%%%%%%%%%%%%%%%%%%%%%%%%%%%%%%%%%%%%%%%%%%%%%%%%%%%%%%%%%%%%%%

%%%%%%%%%%%%%%%% TRẮC NGHIỆM 4 PHƯƠNG ÁN %%%%%%%%%%%%%%%%

% nguồn: MaHoaNH_TracNghiem_Lop10
\begin{ex}%[0D6N1-1]%[ma:92e71e0f]
	Quy tròn số  $2{,}654$ đến hàng phần chục, ta được số $2{,}7$. Sai số tuyệt đối của số gần đúng $2{,}7$ là
	\choice 
		{$-0{,}046$} 
		{\True $0{,}046$} 
		{$2{,}7$} 
		 {$2{,}654$}
	\loigiai{ 
		Sai số  tuyệt đối là $\left|2{,}654-2{,}7\right|=0{,}046$.
	} 
\end{ex}

% nguồn: MaHoaNH_TracNghiem_Lop10
\begin{ex}%[0D6H1-1]%[ma:8bf2da51]
	Cho  giá trị gần đúng của $\dfrac{8}{17}$ là $0{,}47$. Sai số tuyệt đối của $0{,}47$ là
	\choice 
		 {\True $0{,}001$}
		 {$0{,}002$}
		 {$0{,}003$}
		 {$0{,}004$}
	 \loigiai{
		Ta có  $\dfrac{8}{17}=0{,}470588235294 \ldots$.\\
		Sai số tuyệt đối của $0{,}47$ là $\left|0{,}47-\dfrac{8}{17}\right|<|0{,}47-0{,}471|=0{,}001$. 
	} 
\end{ex}

% nguồn: MaHoaNH_TracNghiem_Lop10
\begin{ex}%[0D6H1-1]%[ma:2f320d3b]
	Đo chiều dài của một cây thước, ta được kết quả $\overline{a}=45 \pm 0{,}3$ cm. Khi đó sai số  tuyệt đối của phép đo được ước lượng là
	 \choice
		{$\Delta_{45}=0{,}3$} 
		{\True  $\Delta_{45} \leq 0{,}3$}
		{$\Delta_{45}  \leq -0{,}3$}
		{$\Delta_{45}  = 0{,}3$}
	\loigiai{ 
		Ta có độ dài dài gần đúng của cây thước là $a=45$  với độ chính xác $d=0{,}3$.\\
		Nên sai số tuyệt đối  $\Delta_{45} \leq d=0{,}3$
	 .}
\end{ex}


%%%%%%%%%%%%%%%% ĐÚNG/SAI %%%%%%%%%%%%%%%%

% nguồn: MaHoaNH_DungSai_Lop10
\begin{ex}%[0D6H1-1]%[ma:e6edf17f]
	Bằng phương pháp hình học nhà toán học Acsimet đã  chứng minh được $\dfrac{223}{71} < \pi < \dfrac{22}{7}$.
	Các khẳng  định sau đây đúng hay sai?
	 \choiceTF
		{Sai  số tuyệt đối lớn hơn $0{,}0025$}
		{\True Sai số tuyệt đối không vượt quá $0{,}0022$} 
		{\True Sai số tương đối không  vượt quá $0{,}065\%$}
		{Sai số tương đối  không vượt quá $0{,}0064\%$}
	 \loigiai{
		\begin{itemize}
			\item Ta có $\Delta_{\pi}  = |\overline{\pi}-\pi| < \dfrac{22}{7} - \dfrac{223}{71} \approx 0{,}002$.
			\item Ta có $\delta_{\pi} = \dfrac{\Delta_{\pi}}{\pi}  < \dfrac{0{,}002}{\dfrac{22}{7}} \approx 0{,}065\%$.
		\end{itemize}
	} 
\end{ex}

% nguồn: MaHoaNH_DungSai_Lop10
\begin{ex}%[0D6H1-1]%[ma:9686ac57]
	Hai bạn A và B thực hiện phép đo một cái bàn dài  $\ell=120$ cm. Bạn A có kết quả là $ \ell_1=120{,}4$ cm, bạn B có kết quả là $\ell_2=119{,}6$ cm.
	 Các khẳng định sau đây đúng hay sai?
	 \choiceTF
		{\True  Sai số tuyệt đối phép đo của bạn A là $0{,}4$ cm}
		{Sai số  tuyệt đối phép đo của bạn B là $-0{,}4$ cm}
		{Sai số tương đối phép đo của bạn B  nhỏ hơn $2\%$}
		{\True Sai số tương đối phép đo bạn  A và bạn B bằng nhau}
	\loigiai{ 
		\begin{itemchoice}
			\itemch  Sai số tuyệt đối phép đo của bạn A là $|120,4 -120| = 0{,}4$.
			\itemch Sai số tuyệt  đối phép đo của bạn B là $|119{,}6 -120| = 0{,}4$.
			\itemch Sai số tương đối phép đo của bạn B là $\delta_{\ell_2} = \dfrac{\Delta_{\ell_2}}{|\ell_2|} = \dfrac{0{,}4}{120}  \approx 0{,}333 \%$.
			\itemch  Sai số tương đối phép đo của bạn A là $\delta_{\ell_1} = \dfrac{\Delta_{\ell_1}}{|\ell_1|} = \dfrac{0{,}4}{120} \approx 0{,}333 \%$. \\
			Vậy sai số tương đối phép đo bạn A và bạn  B bằng nhau.
		\end{itemchoice}
	 }
\end{ex}

% nguồn: MaHoaNH_DungSai_Lop10
\begin{ex}%[0D6V1-1]%[ma:c4cc8cf7]
	Kết quả đo chiều dài của một cây cầu được ghi $152$ m $\pm$ $0{,}2$ m; kết quả đo chiều cao  của một ngôi nhà được ghi là $15{,}2$ m $\pm 0{,}1$ m. Khi đó
	 \choiceTF
		{\True Sai số tương đối trong cách  ghi thứ nhất (đo chiều dài của một cây cầu) là $\delta_1\leqslant\dfrac{d_1}{\left|a_1\right|}=\dfrac{0{,}2}{152}\approx 0,13\%$}
		{\True  Sai số tương đối trong cách ghi thứ hai (đo chiều cao của một ngôi nhà) là $\delta_2\leqslant\dfrac{d_2}{\left|a_2\right|}=\dfrac{0{,}1}{15{,}2}\approx 0,66\%$}
		{Sai số tương đối trong cách ghi thứ  hai (đo chiều cao của một ngôi nhà) lớn hơn $0{,}66\%$}
		{Cách ghi thứ nhất (đo chiều dài cây cầu) có độ chính xác thấp hơn cách ghi thứ hai (đo chiều cao ngôi  nhà)}
	\loigiai{ 
		\begin{itemchoice}
			\itemch Đúng. Vì Sai số tương đối trong cách ghi thứ nhất là $\delta_1\leqslant\dfrac{d_1}{\left|a_1\right|}=\dfrac{0{,}2}{152}\approx  0{,}13\%$ (hay sai số tương đối không vượt quá $0{,}13\%)$.
			\itemch Đúng. Vì Sai  số tương đối trong cách ghi thứ hai là $\delta_2\leqslant\dfrac{d_2}{\left|a_2\right|}=\dfrac{0{,}1}{15{,}2}\approx 0{,}66\%$ (hay sai số tương đối không vượt quá $0{,}66\%$).
			\itemch Sai. Vì Qua đánh giá sai số tương đối trong hai cách ghi, ta thấy $0{,}13\%<0{,}66\%$ nên cách ghi thứ nhất (đo chiều dài cây cầu) có độ  chính xác cao hơn cách ghi thứ hai (đo chiều cao ngôi nhà).
			\itemch Sai. Vì Qua đánh giá sai số tương đối trong hai cách ghi, ta thấy $0{,}13\%<0{,}66\%$ nên cách ghi thứ  nhất (đo chiều dài cây cầu) có độ chính xác cao hơn cách ghi thứ hai (đo chiều cao ngôi nhà).
		\end{itemchoice}
	} 
\end{ex}


%%%%%%%%%%%%%%%% TRẢ LỜI NGẮN %%%%%%%%%%%%%%%%

% nguồn: MaHoaNH_TraLoiNgan_Lop10
\begin{bt}%[0D6H1-1]%[ma:6b3b44b9]
	 Làm tròn số $2315{,}56$ đến hàng đơn vị. Tính sai số tuyệt đối của chúng.
	\shortans{0,44} 
	 \loigiai{
		Số $2315{,}56$ được làm tròn đến  hàng đơn vị là $2316$ và sai số tuyệt đối là \\ $\Delta=|2315{,}56-2316|=0{,}44$.
	} 
\end{bt}

% nguồn: MaHoaNH_TraLoiNgan_Lop10
\begin{bt}%[0D6H1-1]%[ma:0a361b25]
	Trong một cuộc điều tra dân số, người ta viết dân số của một tỉnh là $3\,574\,625 \pm 50\,000$ (người). Hãy đánh giá sai số tương đối  của số gần đúng này?(Làm tròn đến hàng phần mười)
	\shortans{1,4} 
	\loigiai{ 
		Ta có $a=3\,574\,625$ và $d=50\,000$ nên sai số tương đối $\delta_{a} \leqslant \dfrac{d}{|a|}=\dfrac{50\,000}{3\,574\,625} \approx 1{,}4  \%$.
	} 
\end{bt}

% nguồn: MaHoaNH_TraLoiNgan_Lop10
\begin{bt}%[0D6H1-1]%[ma:4d3b4a5d]
	Bạn Lan tính diện tích hình tròn bán kính $r=3$ cm bằng công  thức $S=3{,}14 \cdot 3^{2}=28{,}26$ cm$^2$.
	Biết  rằng $3{,}1<\pi<3{,}2$, hãy ước lượng sai số tương đối của $S$. Làm tròn đến hàng phần trăm.
	 \shortans{1,91}
	 \loigiai{
		Diện tích đúng kí hiệu là $\overline{S}$ thỏa mãn $3{,}1\cdot 3^{2}<\overline{S}<3{,}2\cdot 3^{2} \Leftrightarrow  27{,}9<\overline{S}<28{,}8$.\\
		Do đó  $27{,}9-28{,}26<\overline{S}-S<28{,}8-28{,}26 \Leftrightarrow-0{,}36<\overline{S}-S<0{,}54$.\\
		Suy  ra $|\overline{S}-S|<0{,}54$.
		Vậy sai số tương đối không vượt quá $\dfrac{0{,}54}{28{,}26} \approx 1{,}91  \%$.
	} 
\end{bt}

% nguồn: MaHoaNH_TraLoiNgan_Lop10
\begin{bt}%[0D6V1-1]%[ma:80bc1897]
	Độ cao của một ngọn núi đo được  là $h=1\,372{,}5$ m. Với sai số tương đối mắc phải là $0{,}5\%$. Hãy xác định sai số tuyệt đối của kết quả đo trên và viết $h$ dưới dạng chuẩn.
	 \shortans{1,37}
	\loigiai{ 
		Theo công thức $\Delta_h=\dfrac{\Delta_h}{|h|}$ ta có 
		 \[\Delta_h=h \cdot \Delta_h=1\,372 \cdot 5 \cdot \dfrac{0{,}5}{1\,000}=0{,}68\,625.\]
		Và $h$ viết dưới dạng  chuẩn là $h=1\,373$ m.
	 }
\end{bt}

% nguồn: MaHoaNH_TraLoiNgan_Lop10
\begin{bt}%[0D6V1-1]%[ma:e85dde35]
	Một giá trị gần đúng của $\pi$ là $3{,}142$. Hãy ước lượng sai  số tuyệt đối và sai số tương đối của giá trị gần đúng trên biết rằng $3{,}141<\pi<3{,}143$.
	Làm tròn  đến hàng phần trăm.
	\shortans{0,03} 
	\loigiai{ 
		Một giá trị gần đúng của $\pi$  là $3{,}142$. Hãy ước lượng sai số tuyệt đối và sai số tương đối của giá trị gần đúng trên biết rằng $3{,}141<\pi<3{,}143$.\\
		Ta có $3{,}141 - 3{,}142<\pi - 3{,}142<3{,}143  - 3{,}142\Leftrightarrow - 0{,}001<\pi - 3{,}142<0{,}001$.\\ Suy ra$\pi - 3{,}142\mid<0{,}001$.\\
		Do đó sai số tuyệt đối là  $0{,}001$.
		Sai số tương đối không vượt  quá $\dfrac{0{,}001}{|3{,}142|}<\dfrac{0{,}001}{3{,}142}\approx 0{,}0318\%$.
	 }
\end{bt}


%%%%%%%%%%%%%%%% TỰ LUẬN %%%%%%%%%%%%%%%%

% nguồn: MaHoaNH_TuLuan_Lop10
\begin{bt}%[0D6H1-1]%[ma:2b599c86]
	Hãy viết  số quy tròn của số a với độ chính xác $d$ được cho sau đây $\bar{a}=17\,658 \pm 16$.
	 \loigiai{
	Ta có $10<16<100$ nên hàng cao nhất mà $d$ nhỏ hơn một đơn vị của hàng đó là hàng trăm. Do đó ta phải quy tròn số $17\,638$  đến hàng trăm. Vậy số quy tròn là $17\,700$.}
\end{bt}

% nguồn: MaHoaNH_TuLuan_Lop10
\begin{bt}%[0D6H1-1]%[ma:55151986]
	Dùng phân số  $\dfrac{33}{19}$ để làm số gần đúng cho $\sqrt{3}$. Hãy đánh giá sai số tuyệt đối mắc phải là bao nhiêu?
	 \loigiai{
		Sử dụng máy tính cầm tay, ta  có $\sqrt{3}=1{,}732050808\ldots;\dfrac{33}{19}=1{,}736842105\ldots$\\
		Sai số  tuyệt đối của số gần đúng là\\
		$\begin{aligned}
		 \left|\sqrt{3}-\dfrac{33}{19}\right|& =|1,732050808\ldots-1{,}736842105\ldots |\\
		&=1{,}736842105\ldots-1{,}732050808\ldots <1{,}74-1{,}73=0{,}01.\\ 
	\end{aligned}$\\
 }
\end{bt}

% nguồn: MaHoaNH_TuLuan_Lop10
\begin{bt}%[0D6V1-1]%[ma:65f2ff57]
	Bạn Ngân có một mảnh nhựa với bề mặt hình tròn bán kính $1$ dm. Bạn ấy thực hiện đo chu vi của mép mảnh nhựa đó bằng cách sử dụng một sợi dây dài không dãn như sau: Cố định một đầu sợi dây trên mép  mảnh nhựa, rồi quấn sợi dây quanh mép mảnh nhựa một vòng cho đến khi đầu dây cố định chạm vào thân sợi dây lần đầu tiên, sau đó đo độ dài phần dây chạm vào mép mảnh nhựa và được kết quả là $6$ dm. Khi đó sai số tuyệt đối trong phép đo không vượt quá bao nhiêu dm.
	\loigiai{ 
		Chu vi của mép mảnh  nhựa là: $\bar{a}=2\pi \cdot 1=2\pi$ (dm).\\
		Gọi chu vi của mép mảnh nhựa mà bạn Ngân đo được là $a$, suy ra $a=6$ (dm). Vì $3<\pi<3{,}15$ nên $6<2\pi<6{,}3\Rightarrow 6<\bar{a}<6{,}3\Rightarrow\Delta_a=|\bar{a}-6|<6{,}3-6=0,3$.\\ 
	Suy ra sai số tuyệt  đối trong phép đo không vượt quá $0{,}3$ dm.}
\end{bt}
